\documentclass[11pt,a4paper]{article}

\usepackage[utf8]{inputenc}
\usepackage[margin=1in]{geometry}
\usepackage{hyperref}
\usepackage{fancyhdr}
\usepackage{booktabs}
\usepackage{amsmath,amssymb}
\usepackage{float}
\usepackage{enumitem}
\usepackage{lastpage}
\usepackage{array}
\usepackage{multirow}

\hypersetup{colorlinks=true,linkcolor=black,urlcolor=blue,citecolor=black}

\pagestyle{fancy}
\fancyhf{}
\fancyhead[L]{\small Anchor-Based Item Extraction from SEC 10-K Filings}
\fancyhead[R]{\small \thepage/\pageref{LastPage}}
\renewcommand{\headrulewidth}{0.4pt}

\begin{document}

\begin{center}
{\small\scshape Technical Report}\\[2pt]
\rule{\textwidth}{0.5pt}\\[12pt]
{\LARGE\bfseries Anchor-Based Item Extraction from SEC 10-K Filings}\\[12pt]
Eric Bi\\
March 24, 2026\\[8pt]
\end{center}

\vspace{4pt}
\noindent
\begin{tabular}{@{}lll@{}}
\toprule
\textit{Version} & \textit{Date} & \textit{Comments}\\
\midrule
1.0 & March 17, 2026 & Initial release\\
1.1 & March 18, 2026 & Add document-level retrieval rate; DP scoring and shared-anchor improvements\\
1.2 & March 21, 2026 & Improve retrieval 70.9\% $\to$ 77.0\%; TOC page-number fix; EOF extension\\
1.3 & March 24, 2026 & GT noise analysis; adjusted retrieval rate 84.4\%\\
\bottomrule
\end{tabular}

\vspace{12pt}
\begin{center}\textbf{Abstract}\end{center}

\noindent
We present a rule-based pipeline for extracting standardized item sections from SEC 10-K filings. The method exploits the anchor-link structure of the HTML Table of Contents to identify section boundaries, classifies anchors via a three-tier regex hierarchy, and resolves ordering conflicts through a dynamic programming formulation of the weighted longest increasing subsequence problem. On a corpus of 392 annotated filings across three evaluation sets, the pipeline achieves a mean character-level $F_1$ score of \textbf{97.5\%} and a raw document-level retrieval rate of \textbf{77.0\%} (302/392), where a document is ``retrieved'' only when every ground truth item is present, no false positives exist, and all item boundaries exceed $F_1 \geq 0.9$. A ground truth noise analysis reveals that 31 of the 90 remaining failures are attributable to annotation inconsistencies in signatures and item16 labels, yielding a noise-adjusted retrieval rate of \textbf{84.4\%} (331/392).

\vspace{8pt}
\tableofcontents
\newpage

%% =====================================================================
\section{Introduction}
%% =====================================================================

The SEC Form 10-K is a comprehensive annual report required of all publicly traded companies in the United States. Each 10-K follows a prescribed structure of approximately 20 item sections (see Table~\ref{tab:items}), spanning topics from business operations and risk factors to executive compensation and financial statements.

Automated extraction of individual item sections is a prerequisite for downstream tasks such as financial text classification, risk analysis, and regulatory compliance monitoring. However, 10-K filings exhibit substantial formatting heterogeneity: companies use different HTML structures, anchor naming conventions, and section delineation strategies, making robust extraction non-trivial.

This report describes a deterministic, rule-based extraction pipeline that achieves 97.5\% mean character-level $F_1$ and 77.0\% document-level retrieval rate across 392 evaluated filings without reliance on trained models. We detail the method (\S\ref{sec:method}), the evaluation protocol (\S\ref{sec:eval}), the iterative improvement process that raised the retrieval rate from 70.9\% to 77.0\% (\S\ref{sec:iteration}), a ground truth noise analysis establishing an adjusted rate of 84.4\% (\S\ref{sec:noise}), and future directions toward agentic applications (\S\ref{sec:future}).

\begin{table}[H]
\centering
\caption{Standard 10-K item sections.}
\label{tab:items}
\begin{tabular}{ll}
\toprule
\textbf{Item} & \textbf{Description}\\
\midrule
Item 1 & Business overview\\
Item 1A & Risk factors\\
Item 1B & Unresolved staff comments\\
Item 2 & Properties\\
Item 3 & Legal proceedings\\
Item 4 & Mine safety disclosures\\
Item 5 & Market for registrant's common equity\\
Item 6 & Selected financial data / Reserved\\
Item 7 & Management's discussion and analysis\\
Item 7A & Quantitative and qualitative disclosures about market risk\\
Item 8 & Financial statements and supplementary data\\
Item 9 & Changes in and disagreements with accountants\\
Item 9A & Controls and procedures\\
Item 9B & Other information\\
Item 9C & Disclosure regarding foreign jurisdictions\\
Items 10--14 & Directors, executive compensation, security ownership, certain\\
& relationships, principal accountant fees\\
Item 15 & Exhibits and financial statement schedules\\
Item 16 & Form 10-K summary\\
Signatures & Signatory attestation\\
\bottomrule
\end{tabular}
\end{table}


%% =====================================================================
\section{Method}\label{sec:method}
%% =====================================================================

The pipeline consists of five sequential stages. The input is a raw SEC EDGAR submission file (\texttt{.txt}); the output is a dictionary mapping each identified item to its corresponding HTML slice.

\subsection{Stage 1: Document Isolation}

SEC submission files may contain multiple embedded documents (the 10-K proper, exhibits, XBRL data, graphics). We extract the 10-K document by searching for the \texttt{<DOCUMENT>} block whose \texttt{<TYPE>} field matches \texttt{10-K} or \texttt{10-K/A}, then isolate the HTML content within its \texttt{<TEXT>} element.

\subsection{Stage 2: TOC Anchor Collection}

We scan the HTML for all internal hyperlinks of the form \texttt{<a href="\#ID">}. The set of referenced anchor IDs constitutes the candidate pool. This filtering step is essential: a typical 10-K contains hundreds of \texttt{id} attributes for footnotes, cross-references, and formatting; restricting to TOC-referenced anchors reduces the search space by an order of magnitude.

\paragraph{Page-number link handling.} SEC filings use two TOC link patterns. In \emph{descriptive links}, the anchor text contains the item description (e.g., \texttt{<a href="\#id">Item~1A.\ Risk Factors</a>}). In \emph{page-number links}, the item description appears in a separate table cell:

\begin{verbatim}
  <td>Item 1A. Risk Factors</td>
  <td><a href="#id">35</a></td>
\end{verbatim}

\noindent
When the link text is $\leq 6$ characters and matches \verb|\d+|, the parser looks backward in the HTML for contextual item description text using two strategies: (1)~find the nearest preceding \texttt{<tr} tag and classify all text in that table row, or (2)~split on the last \texttt{</a>} tag and classify the remaining text in the same inline context. This handles the common pattern where the item description and page number reside in separate \texttt{<td>} cells.

\subsection{Stage 3: Anchor Element Localization}

For each ID in the candidate pool, we locate the corresponding HTML element by matching \texttt{id} or \texttt{name} attributes across all tag types (\texttt{<a>}, \texttt{<div>}, \texttt{<p>}, \texttt{<span>}, \texttt{<td>}, etc.). When duplicate IDs exist, we retain the \emph{last} occurrence, as the first typically resides in the TOC header while the last marks the actual section body.

\subsection{Stage 4: Anchor Classification}

A typical 10-K filing contains 30--60 TOC-referenced anchors, but only $\sim$20 correspond to actual item sections. The classification task is to determine which anchor corresponds to which 10-K item (or to no item at all).

For each anchor, we extract a \emph{forward lookahead window} and a \emph{backward context window}. These text windows are evaluated against a three-tier classification hierarchy:

\begin{table}[H]
\centering
\caption{Classification tiers and confidence scores.}
\begin{tabular}{rll}
\toprule
\textbf{Conf.} & \textbf{Source} & \textbf{Tier}\\
\midrule
10 & Shared-anchor override (multiple TOC links $\to$ same ID) & 0\\
9 & Anchor ID pattern (e.g., \texttt{ITEM1ARISKFACTORS}) & 0a\\
8 & TOC link text mapping (TOC says ``Item 1A'' $\to$ anchor) & 0b\\
6 & ``Item $X$'' pattern in first 150 chars of lookahead & 1\\
5 & ``Item $X$'' pattern in full lookahead (up to 2000 chars) & 1\\
4 & ``Item $X$'' pattern in backward context & 1\\
3 & Descriptive title in first 150 chars of lookahead & 2\\
2 & Descriptive title in full lookahead & 2\\
1 & Descriptive title in backward context & 2\\
\bottomrule
\end{tabular}
\end{table}

\paragraph{Sequence-constrained selection via dynamic programming.}
Multiple anchors often classify as the same item. We formalize selection as an optimization problem. Let $\mathcal{C} = \{(o_i, s_i, c_i)\}_{i=1}^{N}$ denote the set of candidate anchors, where $o_i$ is the character offset, $s_i \in \{1,2,\ldots,24\}$ is the sequential item index, and $c_i$ is the classification confidence. We seek a subset $\mathcal{A} \subseteq \mathcal{C}$ such that offsets are strictly increasing, item indices are strictly increasing, and $\sum_{i \in \mathcal{A}} w(c_i, o_i)$ is maximized. This is equivalent to the weighted longest increasing subsequence (WLIS) problem, solved via $O(N^2)$ dynamic programming. The weight function $w$ prioritizes classification confidence while preferring \emph{earlier} offsets among same-confidence candidates.

\paragraph{Heading scan fallback.}
After the DP solver produces its assignment, some TOC-mapped items may remain unassigned. For each such item, we scan the document body between the nearest assigned neighbors for an ``Item $X$'' heading pattern. If found at a position consistent with the item sequence, the item is inserted into the assignment.

\subsection{Stage 5: HTML Slicing}

Given the ordered anchor sequence $\{(o_k, \text{item}_k)\}_{k=1}^{K}$, we extract the HTML substring $h[o_k : o_{k+1}]$ for each item $k$. Boundary adjustments handle two edge cases:

\begin{itemize}[nosep]
    \item \textbf{Wrapper \texttt{<div>} elements}: when an anchor is wrapped in \texttt{<div><a id="...">{}</a></div>}, the previous item's slice extends through the closing \texttt{</div>}, and the current item's slice begins at the \texttt{<a>} tag.
    \item \textbf{Terminal item}: the final item extends to \texttt{</body>} or end of the document. An extended variant (\texttt{process\_file\_extended}) extends to end-of-file in the raw submission, matching the ground truth annotation convention.
\end{itemize}

\paragraph{Special cases.}
Signatures suppression: when detected only via a fallback bold-text scan (no TOC link), the entry is suppressed to avoid false positives. Item 16 placeholder detection: when the item contains only ``None'' or ``Not applicable'' and signatures are suppressed, the value is set to empty. Shared-anchor item groups: when multiple TOC entries point to a single anchor (common for Part~III Items 10--14 incorporated by reference), the content is assigned to the last item in the group and empty placeholder keys are emitted for the remaining items.


%% =====================================================================
\section{Evaluation Protocol}\label{sec:eval}
%% =====================================================================

\subsection{Dataset}

The evaluation corpus comprises three sets of 167 SEC EDGAR submission files each. After excluding filings with empty ground truth annotations, the evaluated subset contains 133, 132, and 127 filings for Sets 1, 2, and 3, respectively (392 filings total).

\subsection{Metric: Character-Level $F_1$}

We employ character-level $F_1$ computed over bag-of-character representations. For a predicted string $\hat{y}$ and ground truth string $y$, let $\hat{t} = \text{strip}(\hat{y})$ and $t = \text{strip}(y)$ denote the plain-text representations obtained by removing all HTML tags and decoding all HTML entities. Define character frequency vectors $\boldsymbol{p}, \boldsymbol{g} \in \mathbb{N}^{|\Sigma|}$ over the character alphabet $\Sigma$. Then:
\begin{equation}
P = \frac{\sum_{c \in \Sigma} \min(p_c, g_c)}{\sum_{c \in \Sigma} p_c}, \quad
R = \frac{\sum_{c \in \Sigma} \min(p_c, g_c)}{\sum_{c \in \Sigma} g_c}, \quad
F_1 = \frac{2PR}{P+R}
\end{equation}

\subsection{Metric: Document-Level Retrieval Rate}

Following Zhang et al., we define a stricter document-level metric. A document is considered ``retrieved'' only when \emph{all three} of the following conditions hold simultaneously:
\begin{enumerate}[nosep]
    \item \textbf{No missing items}: every ground truth item key is present in the prediction;
    \item \textbf{No false positives}: no predicted item keys are absent from the ground truth;
    \item \textbf{Correct boundaries}: every item's character-level $F_1 \geq 0.9$.
\end{enumerate}

\noindent
The retrieval rate is the fraction of evaluated documents that satisfy all three conditions.


%% =====================================================================
\section{Results}\label{sec:results}
%% =====================================================================

\subsection{Overall Performance}

\begin{table}[H]
\centering
\caption{Overall performance metrics (after all improvements).}
\label{tab:overall}
\begin{tabular}{lrrrr}
\toprule
\textbf{Dataset} & \textbf{Files} & $\boldsymbol{F_1}$ & \textbf{Extraction Rate} & \textbf{Doc Retrieval}\\
\midrule
Set 1 & 133 & 98.0\% & 98.8\% & 81.2\%\\
Set 2 & 132 & 97.5\% & 98.4\% & 73.5\%\\
Set 3 & 127 & 97.3\% & 97.9\% & 76.4\%\\
\midrule
\textbf{Mean} & \textbf{392} & \textbf{97.6\%} & \textbf{98.4\%} & \textbf{77.0\%}\\
\bottomrule
\end{tabular}
\end{table}

\subsection{Error Analysis}

Of the 90 documents that fail the retrieval criterion, the failure modes decompose as follows:

\begin{itemize}[nosep]
    \item \textbf{Boundary errors only} (40 docs): All items are present with no false positives, but at least one item has $F_1 < 0.9$. The most affected items are signatures (20 docs, where ground truth extends into exhibit content beyond the 10-K document boundary) and item 16 (14 docs).
    \item \textbf{Missing items only} (7 docs): Dominated by signatures (5 docs), where the filing lacks a TOC entry for signatures and the fallback detector is suppressed to avoid false positives.
    \item \textbf{False positives only} (11 docs): Extra predicted items not in ground truth, primarily item 16 (7 docs) in filings where the ground truth omits this optional section.
    \item \textbf{Mixed failures} (32 docs): Combinations of the above, often involving filings with non-standard HTML structure.
\end{itemize}


%% =====================================================================
\section{Retrieval Rate Iteration}\label{sec:iteration}
%% =====================================================================

Starting from a baseline retrieval rate of 70.9\% (278/392), we conducted a systematic improvement effort. This section documents the two changes that produced gains and the ten approaches that were reverted.

\subsection{Root Cause Discovery: Page-Number TOC Links}

The primary root cause was identified through reverse-engineering the ground truth annotations. The GT tool and our pipeline use the \emph{same} set of TOC-referenced anchors. The key gap was in TOC link classification: many filings encode page numbers as separate link elements in a different table cell from the item description (\S2.2). Our baseline parser classified only the link text, missing these anchors entirely.

\subsection{Fix 1: Page-Number TOC Link Classification (+22 docs)}

We extended \texttt{parse\_toc\_links} with the table-row and same-cell context strategies described in \S2.2. This fix is applied during TOC parsing, before anchor classification. It does not change the classification tiers or the DP solver---it simply allows more anchors to receive correct item labels from the TOC.

\subsection{Fix 2: Last-Item End-of-File Extension (+2 docs)}

We implemented \texttt{process\_file\_extended()}, which extends the last extracted item to the end of the raw submission file. A helper function \texttt{\_extract\_10k\_full()} returns both the 10-K HTML text and the full file content with offset mapping.

\subsection{Cumulative Impact}

\begin{table}[H]
\centering
\caption{Document-level retrieval rate before and after improvements.}
\begin{tabular}{lrrr}
\toprule
\textbf{Set} & \textbf{Baseline} & \textbf{After fixes} & \textbf{Change}\\
\midrule
Set 1 & 99/133 (74.4\%) & 108/133 (81.2\%) & +9\\
Set 2 & 90/132 (68.2\%) & 97/132 (73.5\%) & +7\\
Set 3 & 89/127 (70.1\%) & 97/127 (76.4\%) & +8\\
\midrule
\textbf{Overall} & \textbf{278/392 (70.9\%)} & \textbf{302/392 (77.0\%)} & \textbf{+24}\\
\bottomrule
\end{tabular}
\end{table}

\subsection{Negative Results}

Ten approaches were implemented, evaluated, and reverted for causing net regressions. Table~\ref{tab:experiments} provides the complete experiment log. The most instructive failures include: trailing marker stripping ($-21$ percentage points; GT \emph{includes} structural markers as part of slices), blanket signatures un-suppression (26 fixes but 75 new false positives), and ML signatures prediction (92\% accuracy identical to the heuristic baseline; does not generalize across evaluation sets).

\begin{table}[H]
\centering
\caption{Complete experiment log, ordered by net impact.}
\label{tab:experiments}
\begin{tabular}{>{\raggedright}p{5.8cm}rrl}
\toprule
\textbf{Approach} & \textbf{Before} & \textbf{After} & \textbf{Status}\\
\midrule
Page-number TOC link classification & 278 & 300 & \textbf{Kept}\\
Last-item EOF extension & 300 & 302 & \textbf{Kept}\\
First-body anchor preference & 302 & 302 & Neutral\\
DP solver: prefer later positions & 302 & 295 & Reverted\\
Extended heading scan (non-TOC) & 300 & 298 & Reverted\\
Expanded title patterns & 300 & 297 & Reverted\\
Unicode normalization & 300 & 298 & Reverted\\
Item16 trailing trim & 278 & 277 & Reverted\\
ML signatures prediction & 278 & 270 & Reverted\\
ML post-processing pipeline & 278 & 244 & Reverted\\
Trailing marker stripping & 278 & 219 & Reverted\\
Blanket signatures un-suppression & 278 & 229 & Reverted\\
\bottomrule
\end{tabular}
\end{table}


%% =====================================================================
\section{Ground Truth Noise Analysis}\label{sec:noise}
%% =====================================================================

Inspection of failing documents reveals that the ground truth annotations contain systematic noise, particularly for signatures and item 16. We classify each item-level failure as either \emph{GT noise} (annotation inconsistency) or \emph{real pipeline error} based on objective criteria.

\subsection{Noise Classification Criteria}

\begin{itemize}[nosep]
    \item \textbf{Boundary signatures --- source file mismatch}: GT value exceeds 1\,M characters (12--86\,MB), indicating that GT was generated from a different version of the source file with different HTML encoding. Character-level divergence prevents $F_1 \geq 0.9$ even when anchor positions are identical.
    \item \textbf{Boundary item16 --- source file mismatch}: Same pattern as signatures; GT exceeds 1\,M characters or both GT and prediction are short placeholders ($<$100 chars text) with minor HTML differences.
    \item \textbf{Missing signatures --- GT empty/trivial}: GT includes a \texttt{signatures} key with empty or trivial content ($<$100 chars). The pipeline detects signatures but suppresses it to avoid false positives; GT inconsistently includes the key.
    \item \textbf{FP item16 --- GT inconsistently omits}: Pipeline correctly extracts item 16 (present in TOC with placeholder content), but GT annotation omits the key for this filing while including it for other filings with identical content.
    \item \textbf{FP signatures --- GT inconsistently omits}: Filing contains a clear SIGNATURES heading detected by the pipeline, but GT does not include the key.
\end{itemize}

\subsection{Results}

\begin{table}[H]
\centering
\caption{Failure classification and adjusted retrieval rates.}
\label{tab:noise}
\begin{tabular}{lr}
\toprule
\textbf{Category} & \textbf{Documents}\\
\midrule
Currently passing & 302\\
Failing --- all failures are GT noise & 31\\
Failing --- mixed (noise + real errors) & 14\\
Failing --- all failures are real errors & 47\\
\midrule
Total evaluated & 392\\
\bottomrule
\end{tabular}
\end{table}

\begin{table}[H]
\centering
\caption{Retrieval rate under different GT quality assumptions.}
\begin{tabular}{llr}
\toprule
\textbf{Metric} & \textbf{Computation} & \textbf{Rate}\\
\midrule
Raw retrieval rate & 302 / 392 & 77.0\%\\
Noise-adjusted (noise docs $=$ pass) & 331 / 392 & \textbf{84.4\%}\\
Clean GT only (exclude all noise docs) & 302 / 347 & 86.5\%\\
\bottomrule
\end{tabular}
\end{table}

\noindent
The 31 all-noise documents break down as: boundary signatures from source file mismatch (13), boundary item16 from source file mismatch or trivial difference (8), missing signatures with empty GT (5), FP item16 with GT inconsistently omitting (4), and FP/missing signatures (1).

Among the 47 real-failure documents, no single item dominates---failures are distributed across all items (item7: 6 missing, item4: 6 boundary, item1: 5 boundary, etc.), confirming that the pipeline handles signatures and item16 correctly and that the apparent failures on those items are GT noise.


%% =====================================================================
\section{Future Directions: Agentic Applications}\label{sec:future}
%% =====================================================================

The extraction pipeline produces structured, per-item representations of 10-K filings at 97.6\% fidelity with 77.0\% document-level retrieval (84.4\% noise-adjusted). This structured output serves as a foundation for autonomous agent systems that perform financial analysis tasks requiring targeted access to specific filing sections.

\paragraph{Multi-agent risk assessment.} A system of cooperating agents could consume Item~1A (Risk Factors) slices, apply named entity recognition and topic modeling, compare risk taxonomies across consecutive fiscal years, and aggregate findings across an industry sector to identify systemic risks.

\paragraph{Autonomous due diligence.} An agent tasked with pre-acquisition due diligence could extract Items~1, 3, 7, and 8 for a target company's last five fiscal years, identify litigation trends, correlate management outlook with quantitative data, and produce a structured due diligence memo with flagged anomalies.

\paragraph{Regulatory compliance monitoring.} A continuously running agent could poll EDGAR for new filings, run the extraction pipeline, and compare extracted Items~9A and 9C against a configurable rule set encoding regulatory requirements, flagging filings that omit required disclosures.

\paragraph{Conversational financial analyst.} An LLM-based agent could use the pipeline as a tool within a retrieval-augmented generation (RAG) architecture, applying an \emph{extract-then-retrieve} pattern that avoids the noise and irrelevance of embedding and searching entire filings.

\paragraph{Portfolio surveillance swarm.} A swarm of lightweight agents (one per portfolio holding) could each track its assigned company's EDGAR filings, extract all items, compute diffs against prior years, and route anomaly reports to a portfolio manager agent for daily surveillance briefings.


%% =====================================================================
\section{Conclusion}
%% =====================================================================

We presented a deterministic, rule-based pipeline for extracting standardized item sections from SEC 10-K filings. The method exploits the anchor-link structure of HTML Tables of Contents, classifies anchors via a three-tier regex hierarchy, and resolves ordering constraints via dynamic programming.

On 392 evaluated filings, the pipeline achieves a mean character-level $F_1$ of 97.6\% and a raw document-level retrieval rate of 77.0\% (302/392), improved from 70.9\% through a targeted fix to TOC page-number link classification. A ground truth noise analysis reveals that 31 of the 90 remaining failures stem from annotation inconsistencies---primarily source file mismatches for signatures and item16 boundaries---yielding a noise-adjusted retrieval rate of \textbf{84.4\%} (331/392).

The extensive negative-result catalog (10 reverted approaches) demonstrates that the pipeline operates near a local optimum for rule-based methods. The high-fidelity, per-item structured output serves as an enabling primitive for agentic financial analysis systems, from multi-agent risk assessment to autonomous due diligence and portfolio surveillance.

\end{document}
